\subsection{亚连续双线性映射}

在下面，我们将定义一个介于连续双线性映射的概念与分别连续双线性映射的概念之间的概念。

命题 3. --- 设 E、F、G 是三个局部凸空间，S 是 E 的有界子集的一个族。设 u 是从 E × F 到 G 中的一个分别连续双线性映射。下列性质等价：

\begin{enumerate}
\def\labelenumi{\alph{enumi})}
\item
  对 0 的每个邻域 \(\mathbf{W}\)（在 \(\mathbf{G}\) 中）和每个集 \(\mathbf{M} \in \mathfrak{S}\)，存在 0 的一个邻域 \(\mathbf{V}\)（在 \(\mathbf{F}\) 中），使得 \(u(\mathbf{M} \times \mathbf{V}) \subset \mathbf{W}\)。
\item
  对每个集 \(\mathbf{M} \in \mathfrak{S}\)，\(\mathbf{M}\) 在映射 \(x \mapsto u(x, .)\) 下的像是 \(\mathcal{L}(\mathrm{F}; \mathrm{G})\) 的一个等度连续子集。
\item
  映射 \(y \mapsto u(., y)\)（从 \(\mathbf{F}\) 到 \(\mathcal{L}_{\mathfrak{S}}(\mathbf{E}; \mathbf{G})\) 中）是连续的。
\item
  表明 \(y \mapsto u(., y)\) 在点 0 处连续，这是根据 \(\mathcal{L}_{\mathfrak{S}}(\mathrm{E}; \mathrm{G})\) 中 0 的邻域的定义（III, 第 13 页）；类似地，a) 表明 M 在映射 \(x \mapsto u(x, .)\) 下的像在点 0 处等度连续（III, 第 16 页）。
\end{enumerate}

定义 2. --- 设 u 是从 E × F 到 G 中的一个双线性映射。我们称 u 为 S-亚连续的，如果 u 分别连续且它满足命题 3 的等价条件 a)、b)、c) 之一。

命题 3 的条件 \(c)\) 表明，\(\mathfrak{S}\) -亚连续双线性映射的概念对 \(\mathfrak{S}\) 的依赖仅通过 \(\mathfrak{S}\) -拓扑（在 \(\mathcal{L}(\mathrm{E},\mathrm{G})\) 上的）实现。

对 F 的有界子集的每个集 T，我们类似地定义 T-亚连续映射的概念，只要在命题 3 中交换 E 与 F 的角色。称一个分别连续双线性映射 u 为 \((\mathfrak{S}, \mathfrak{T})\) -亚连续的，如果它既是 S-亚连续的又是 T-亚连续的。

每个从 \(E \times F\) 到 G 中的连续双线性映射对有界子集的集的每个对 \((\mathfrak{S}, \mathfrak{T})\) 都是 \((\mathfrak{S}, \mathfrak{T})\) -亚连续的：对 G 中 0 的每个邻域 W，存在 E 中 0 的一个邻域 U 和 F 中 0 的一个邻域 V，使得 \(u(U \times V) \subset W\)；由于每个集 \(M \in S\) 有界，存在 \(\lambda > 0\) 使得 \(\lambda M \subset V\)，于是

\[
u (\mathbf {M} \times \lambda \mathbf {V}) = u (\lambda \mathbf {M} \times \mathbf {V}) \subset u (\mathbf {U} \times \mathbf {V}) \subset \mathbf {W}.
\]

逆命题一般不成立（III, 第 47 页, 习题 3）。

命题 4. --- 设 u 是从 E × F 到 G 中的一个 S-亚连续双线性映射。对每个集 M ∈ S，u 在 M × F 上的限制是连续的，且对 F 的每个有界子集 Q，u(M × Q) 在 G 中有界。

第一个断言由 GT, X, § 2, No.~1 的推论 3 得出。设 W 是 G 中 0 的一个邻域；由假设，存在 F 中 0 的一个邻域 V，使得 \(u(\mathbf{M} \times \mathbf{V}) \subset \mathbf{W}\)。由于存在 \(\lambda \neq 0\) 使得 \(\lambda \mathbf{Q} \subset \mathbf{V}\)，我们有 \(\lambda u(\mathbf{M} \times \mathbf{Q}) = u(\mathbf{M} \times \lambda \mathbf{Q}) \subset \mathbf{W}\)，这就证明了命题的第二部分。

命题 5. --- 设 u 是从 E × F 到 G 中的一个 (S, T)-亚连续双线性映射。对每对集 M ∈ S、N ∈ T，u 在 M × N 上一致连续。命题由 GT, X, § 2, No.~1 的命题 2 和 GT, X, § 2, No.~2 的命题 5 立即得出。

命题 6. --- 若 F 是桶型空间，则每个从 E × F 到局部凸空间 G 中的分别连续双线性映射 u 对 E 的有界子集的每个集 S 都是 S-亚连续的。

换言之，线性映射 \(y \mapsto u(., y)\)（从 \(F\) 到 \(\mathcal{L}_b(E; G)\) 中）是连续的。只需（III, 第 30 页, 命题 3）证明：每个有界子集 \(M\)（\(E\) 的）在 \(x \mapsto u(x, .)\) 下的像在 \(\mathcal{L}(F; G)\) 中等度连续。但由命题 1（III, 第 28 页），这个像是 \(\mathcal{L}(F; G)\) 的一个简单有界子集，而由于 \(F\) 是桶型的，\(\mathcal{L}(F; G)\) 的每个简单有界子集都是等度连续的（III, 第 25 页, 定理 1）。

评注. --- 假设 F 的拓扑是 F 上使线性映射 \(h_{\alpha}:F_{\alpha}\to F\) 连续的最细局部凸拓扑（II, 第 27 页）。则命题 3 的条件 c)（III, 第 30 页）表明：若 E 和 G 是局部凸的，则双线性映射 \(u:E\times F\to G\) 是 S-亚连续的，当且仅当诸双线性映射

\[
(x, y _ {\alpha}) \mapsto u (x, h _ {\alpha} (y _ {\alpha}))
\]

（从 \(\mathbf{E} \times \mathbf{F}_{\alpha}\) 到 \(\mathbf{G}\) 中）中的每一个都是 \(\mathfrak{S}\) -亚连续的。

现在假设 \(\mathbf{E}\) 是一个局部凸空间，它是闭向量子空间的一个递增序列 \((\mathbf{E}_n)\)（\(\mathbf{E}\) 的）的严格归纳极限（II, 第 33 页）；则每个集 \(\mathbf{M} \in \mathfrak{S}\) 含于诸 \(\mathbf{E}_n\) 之一且在这个子空间中有界（III, 第 5 页, 命题 6）。我们用 \(\mathfrak{S}_n\) 记属于 \(\mathfrak{S}\) 且含于 \(\mathbf{E}_n\) 的所有子集所成的族。

命题 3 的条件 a)（III, 第 30 页）表明：为使双线性映射 \(u: \mathrm{E} \times \mathrm{F} \to \mathrm{G}\) 是 \(\mathfrak{S}\) -亚连续的，必要且充分的是诸限制 \(u_n: \mathrm{E}_n \times \mathrm{F} \to \mathrm{G}\)（\(u\) 的）中的每一个都是 \(\mathfrak{S}_n\) -亚连续的。

